% ============================================================================
% Supplementary Information SI-2: Generative-Loop Reaction Discovery
%
% Drop-in subsection fragment. Parent SI must provide: amsmath/amssymb,
% booktabs, siunitx (with \angstrom). Bibliography handled at the parent
% level. (Parameter reference lives in si_parameters.tex.)
% ============================================================================

\subsection{Generative-Loop Reaction Discovery}
\label{si:generative}

Each iteration processes a batch of $32$ \textbf{exploration contexts} (a reactant geometry --- single molecule or merged pair --- plus the artefacts the pipeline attaches downstream: TS conformer, product conformer, IRC trajectories, fragment list, barriers). The pipeline has eight stages: (1)~sampling, (2)~merge enhancement, (3)~MoreRed diffusion, (4)~forward barrier gate, (5)~Hessian + IRC validation, (6)~endpoint identity, (7)~fragmentation detection, (8)~backward barrier gate. Critic stages (4--8) reuse the shared machinery in SI-\ref{si:shared} (ML force field, geometry primitives, optimisers); this section documents what is specific to the generative path.


%-----------------------------------------------------------------------------
\subsubsection{Stage 1 --- Concentration-weighted context sampling}
\label{sec:sampling}
%-----------------------------------------------------------------------------

Half of each batch is single-molecule, half is bimolecular merge pairs. The eligible compound pool is scoped to the current experiment and excludes out-of-scope frontier compounds.

\paragraph{Single-molecule sampling.}
A species is drawn from the kinetic snapshot's steady-state distribution (SI-\ref{si:kinetics}); compounds with $n_\text{atoms} < 3$ are excluded. Within the chosen compound, the conformer is drawn from the room-temperature Boltzmann distribution,
\begin{equation}
    P(m) = \frac{\exp(-(E_m - E_\text{min}) / k_B T)}{\sum_{m'} \exp(-(E_{m'} - E_\text{min}) / k_B T)}, \qquad k_B T = \SI{0.0257}{eV}.
\end{equation}

\paragraph{Bimolecular pair sampling.}
$50 \times n_\text{needed}$ candidate pairs are drawn from the steady-state distribution, then filtered against (i)~a composition cap of $\{C{:}5,\ H{:}10,\ O{:}5\}$ (tetrose-sized merged complexes) and (ii)~a per-pair coverage cap of $4$ non-manual reactions in the current experiment. For each surviving pair, one minimum is drawn from each compound at random. Duplicate pairs are kept --- each becomes an independent rotation/noise seed downstream.

\paragraph{Initial geometric merge.}
Each component receives an independent random $SO(3)$ rotation; a random unit vector defines the separation direction; the smallest displacement that yields a non-negative minimum van der Waals gap $\min_{ij}\bigl(\lVert r_i - r_j \rVert - r^{\text{vdW}}_i - r^{\text{vdW}}_j\bigr)$ is found by binary search, then $\SI{0.2}{\angstrom}$ of repulsive buffer is added. Atoms are interleaved by stable-sorting on atomic number; the merged structure is COM-centred. Atom ordering is fully deterministic so that downstream IDs match across workers.

\noindent\begin{tabular}{lr}
\toprule
Parameter & Value \\
\midrule
Batch size & 32 contexts \\
Single/merge split & 50/50 \\
Oversample factor (merges) & 50 \\
Min compound size & 3 atoms \\
Composition cap (merge) & $\{C{:}5,H{:}10,O{:}5\}$ \\
Per-pair reaction cap & 4 \\
$k_B T$ (conformer pick) & \SI{0.0257}{eV} \\
vdW separation buffer & \SI{0.2}{\angstrom} \\
\bottomrule
\end{tabular}

The 50/50 split is hardcoded because a uniform pair draw with $n$ compounds spends $\sim 1/n$ on each cross-pair, making bimolecular discovery rare without a fixed budget. The tetrose composition cap is the dominant top-down constraint on system reach; relaxing it is the principal scale-up knob for chemistry beyond formose-style sugars.


%-----------------------------------------------------------------------------
\subsubsection{Stage 2 --- Merge encounter enhancement}
\label{sec:merge-prep}
%-----------------------------------------------------------------------------

The random encounter geometry from Stage 1 is geometrically valid but typically far from any saddle. Pre-conditioning reduces the burden on the diffusion stage. For each merge context: (i)~generate $5$ additional random merge geometries and pick the lowest-ML-energy candidate; (ii)~run $5$ FIRE micro-relaxation steps with $f_\text{max} = \SI{0.05}{eV/\angstrom}$ to settle the fragments into a reasonable encounter geometry without converging to the encounter-complex minimum; (iii)~reject the context if the binding energy $E_\text{bind} = E(A) + E(B) - E(\text{complex}) > \SI{0.3}{eV}$ (deep encounter complexes are kinetic traps from which both IRC directions relax back into the complex); (iv)~apply per-atom Gaussian noise of $\sigma = \SI{0.15}{\angstrom}$ to break symmetry, since MoreRed is equivariant and produces no movement on a symmetric stationary input; (v)~re-centre for the translation-invariant diffusion.

\noindent\begin{tabular}{lr}
\toprule
Parameter & Value \\
\midrule
Extra orientations ($K$) & 5 \\
FIRE micro-relax steps & 5 \\
FIRE $f_\text{max}$ & \SI{0.05}{eV/\angstrom} \\
Symmetry-break $\sigma$ & \SI{0.15}{\angstrom} \\
Encounter-complex cap & \SI{0.3}{eV} \\
Min pairwise distance (overlap guard) & \SI{0.8}{\angstrom} \\
\bottomrule
\end{tabular}

The \SI{0.3}{eV} binding cap is a chemistry choice: physisorbed and H-bonded complexes are kept, while covalently bound encounter complexes (themselves minima) are rejected --- their reactions out are PES-loop business (SI-\ref{si:pes}). The \SI{0.15}{\angstrom} displacement is calibrated against the diffusion-model noise scale: at $t \sim 200$ on the polynomial schedule the cumulative variance is comparable, so the denoiser sees a perturbation it can act on without losing the input identity.


%-----------------------------------------------------------------------------
\subsubsection{Stage 3 --- Diffusion-based TS proposer}
\label{sec:diffusion}
%-----------------------------------------------------------------------------

\paragraph{Model.}
A SchNetPack equivariant graph network (Sch{\"u}tt et al.\ 2023), trained as a MoreRed-JT joint-time denoiser (Kahouli et al.\ 2024) on relaxed transition states. Per atom it predicts a noise vector $\boldsymbol{\epsilon}_\theta(\mathbf{x}_t, t) \in \mathbb{R}^3$ and a time-step $\hat{t}_\theta(\mathbf{x}_t)$ used for adaptive convergence. The forward process is a variance-preserving Gaussian DDPM (Ho et al.\ 2020) of length $T = 1000$ with a polynomial noise schedule (Hoogeboom et al.\ 2022),
\begin{equation}
    \bar{\alpha}_t = (1 - s)\bigl(1 - (t/T)^p\bigr)^2 + s, \qquad p = 2,\ s = 10^{-5},
\end{equation}
clipped so $\alpha_t / \alpha_{t-1} \geq 10^{-3}$, with forward conditional $q(\mathbf{x}_t \mid \mathbf{x}_0) = \mathcal{N}\bigl(\sqrt{\bar{\alpha}_t}\,\mathbf{x}_0,\,(1 - \bar{\alpha}_t)\,\mathbf{I}\bigr)$ and the standard DDPM lower-bound posterior variance for the reverse step. Translation invariance is enforced by re-centring at every step.

\paragraph{Partial-noise inference.}
Unlike text-to-image diffusion, the TS proposer does not sample from pure Gaussian noise: it partially noises a chemically meaningful input and then denoises. Each batch samples a single noise level $t \sim \mathcal{U}[100, 650]$, forward-diffuses each starting geometry as $\mathbf{x}_t = \sqrt{\bar{\alpha}_t}\,\mathbf{x}_0 + \sqrt{1 - \bar{\alpha}_t}\,\boldsymbol{\epsilon}$ with $\boldsymbol{\epsilon} \sim \mathcal{N}(0, I)$, and reverse-diffuses with up to $1000$ steps:
\begin{equation}
    \mathbf{x}_{t-1} = \tfrac{1}{\sqrt{\alpha_t}}\bigl(\mathbf{x}_t - \tfrac{1-\alpha_t}{\sqrt{1 - \bar{\alpha}_t}}\,\boldsymbol{\epsilon}_\theta\bigr) + \sigma_t\,\mathbf{z}, \quad \mathbf{z} \sim \mathcal{N}(0,I).
\end{equation}
Because the model predicts a per-molecule time-step, the reverse loop is adaptive: molecules whose internal clock decays faster converge in fewer steps, while harder cases use more. Once a molecule's predicted time reaches the convergence threshold ($0$) it is held fixed while the rest continue.

\paragraph{Outputs.}
Each context receives a denoised \texttt{ts\_conformer} (COM-centred), an ML \texttt{ts\_energy}, a geometry hash for cross-worker deduplication, and discovery metadata (\texttt{discovery\_method = "generative"}, the injected noise level $t$, timestamp).

\noindent\begin{tabular}{lr}
\toprule
Parameter & Value \\
\midrule
$T$ (schedule length) & 1000 \\
Schedule shape & polynomial, $p=2$, $s = 10^{-5}$ \\
Clip floor & $10^{-3}$ \\
Variance type & lower-bound \\
Partial-noise range & $[100, 650]$ \\
Max reverse steps & 1000 \\
Convergence step & 0 \\
Cutoff (neighbour list) & \SI{5.0}{\angstrom} \\
\bottomrule
\end{tabular}

The $[100, 650]$ noise range is the central hyperparameter: too low and the denoiser barely moves the input (most candidates are reactant copies); too high and the input is washed out (chemistry is lost). The chosen range keeps $\bar{\alpha}_t \in [0.86, 0.005]$, the regime where the noised input still encodes meaningful structure but the denoiser has freedom to move atoms across bonds. The injected $t$ and the model-predicted $\hat{t}_\theta$ are not constrained to agree --- the reverse loop adapts to the noise level it actually perceives, which gives robustness to off-distribution inputs (e.g., symmetry-broken encounter complexes from Stage 2).


%-----------------------------------------------------------------------------
\subsubsection{Stage 4 --- Forward barrier gate}
\label{sec:fwd-barrier}
%-----------------------------------------------------------------------------

A cheap ML-energy filter. The forward barrier is
\begin{equation}
    \Delta E_\text{fwd} =
    \begin{cases}
        E_\text{TS} - (E_A + E_B) & \text{merged, components } A,B,\\
        E_\text{TS} - E_\text{start}  & \text{single-molecule,}
    \end{cases}
\end{equation}
accepted iff $\Delta E_\text{fwd} \in [\SI{0}{eV}, \SI{2.7211}{eV}]$ (the upper bound is $0.1$~Hartree).

The \SI{2.7211}{eV} ceiling is deliberately loose. At \SI{300}{K} the Eyring prefactor $k_B T/h \approx 6 \times 10^{12}\ \text{s}^{-1}$, so even \SI{2.7}{eV} corresponds to a unimolecular rate $\sim 10^{-33}\ \text{s}^{-1}$ --- kinetic relevance is enforced downstream by the ODE solver, not here. Tightening this cap would risk discarding reactions whose ML barrier is overestimated relative to DFT.


%-----------------------------------------------------------------------------
\subsubsection{Stage 5 --- Hessian and IRC validation}
\label{sec:hessian-irc}
%-----------------------------------------------------------------------------

The bidirectional IRC algorithm (Hessian projection, adaptive displacement, bidirectional Newton relax, energy-based direction assignment) is shared with the PES loop and described in \S\ref{sec:shared-irc}. Generative-specific choices:

\paragraph{Hessian acceptance.}
The candidate is accepted as a TS iff \emph{at least one} significant eigenvalue ($|\lambda_i| > 10^{-4}\,\text{eV/\AA}^2$; \S\ref{sec:shared-hessian}) satisfies $\lambda_i < \SI{-0.01}{eV/\angstrom^2}$. The eigenvector of the \emph{most} negative eigenvalue is taken as the imaginary mode, and this single mode is followed by the bidirectional IRC. We deliberately do \emph{not} require the strict ``exactly one negative eigenvalue'' criterion that the PES loop enforces (\S\ref{sec:pes-prfo}): higher-order signatures are common at diffusion-proposed geometries and forcing strict first-order acceptance would tank recall.

\emph{Consequences of admitting higher-order saddles.} Once the candidate passes this gate, no downstream stage re-counts the negative eigenvalues. A geometry with $\geq 2$ negative modes therefore enters the pipeline if the dominant mode is selected; the subsequent IRC, endpoint, and barrier gates filter such candidates only \emph{implicitly}: a true higher-order saddle is statistically unlikely to (i) relax cleanly to two distinct minima within $350$ Newton steps on each side, (ii) have its dominant mode connect two endpoints that satisfy the endpoint-identity test (Stage 6), and (iii) yield both a forward and a backward barrier inside $(0, \SI{2.7211}{eV}]$ (Stages 4, 8). Pathological higher-order candidates typically fail one of these implicit filters --- most often the bidirectional-barrier check, where the secondary negative mode causes one endpoint to land at an energy above the TS. But candidates whose dominant mode is a sensible reaction coordinate \emph{can} survive even when secondary negative modes are present; we accept that this means the generative loop occasionally admits geometries that are not strictly first-order saddles. Transition-state theory assumes a first-order saddle for the Eyring prefactor, so these admissions carry a small kinetic-rate bias, mitigated downstream by the PBE0/def2-TZVPP barrier refinement (\S\ref{sec:kinetics-dft}) which is computed at the same in-box geometry and re-anchors the rate constants on DFT energetics.

\paragraph{Adaptive displacement constants.}
The clip range is wider than the PES loop's because diffusion-proposed saddles span a broader curvature distribution:
\begin{equation}
    \Delta x = \mathrm{clip}\bigl(\tfrac{0.2}{\sqrt{|\lambda_\text{TS}|}},\ \SI{0.05}{\angstrom},\ \SI{0.5}{\angstrom}\bigr).
\end{equation}

\paragraph{Endpoint relaxation budget.}
Trust-region Newton (\S\ref{sec:shared-newton}) with max $350$ steps (longer than the PES loop's $120$, again reflecting the broader candidate distribution), $f_\text{max} = \SI{0.005}{eV/\angstrom}$, Hessian re-traced every $5$ steps. Acceptance: convergence or bail-late at $f_\text{max} < 10\times$ tolerance.

\paragraph{Context rebuild after direction flip.}
When the IRC reactant (energy-determined; \S\ref{sec:shared-irc}) differs from the sampler's source compound, the context is rebuilt against the actual reactant: re-fragment, re-derive SMILES, re-relax each fragment on its own PES, re-derive merge/single status. This is what lets the generative loop discover reactions whose reactant the sampler did not anticipate.

\noindent\begin{tabular}{lr}
\toprule
Parameter & Value \\
\midrule
Eigenvalue threshold (TS) & \SI{-0.01}{eV/\angstrom^2} \\
Significance mask & $10^{-4}\,\text{eV/\AA}^2$ \\
Adaptive step formula & $0.2 \cdot |\lambda|^{-1/2}$~\AA, clamped $[0.05, 0.5]$ \\
IRC max steps / direction & 350 \\
IRC $f_\text{max}$ & \SI{0.005}{eV/\angstrom} \\
Bail-late tolerance & $10\times f_\text{max}$ \\
\bottomrule
\end{tabular}

Energy-based direction assignment is non-trivial: without it, the same reaction discovered from $A$ vs.\ $B$ of an $A \rightleftharpoons B$ pair would enter the graph in opposite orientations, and downstream dedup would have to handle the symmetric case explicitly. The bail-late tolerance is generous relative to the convergence criterion but appropriate at ML level; the DFT refinement (\S\ref{sec:kinetics-dft}) supersedes the ML barriers consumed by the kinetics solver downstream.


%-----------------------------------------------------------------------------
\subsubsection{Stage 6 --- Endpoint identity}
\label{sec:endpoints}
%-----------------------------------------------------------------------------

\paragraph{Strict mode.}
Exactly one endpoint must match the known reactant, with the match defined by permutation-invariant RMSD (\S\ref{sec:shared-rmsd}) below $\SI{1.0}{\angstrom}$. For single-molecule contexts the reactant is the source compound; for merge contexts the endpoint is fragmented (\S\ref{sec:shared-fragments}) and the validator tries both fragment-to-component assignments. Three outcomes: \textsc{one match} (accept), \textsc{both match} (no new product), \textsc{neither match} (fall through to greedy).

\paragraph{Greedy fallback.}
Enabled by default for both merge and single contexts. The validator drops the source-compound requirement and accepts the reaction iff the two endpoints describe chemically different states: each side fragmented into $\leq 2$ components with canonically distinct SMILES (\S\ref{sec:shared-smiles}) multisets, and no \emph{spectator species} (a fragment that appears on both sides at identical multiplicity while the other species differs --- the reaction $A + B \to A + C$ is rejected because $A$ is unchanged, but $2A \to A + B$ is accepted because $\{A,A\} \neq \{A,B\}$).

Without greedy mode, $> 80\%$ of valid TS proposals would be rejected because the IRC found a reaction whose reactant differs from the sampler's pick (different conformer, tautomer, or related-but-distinct species). Greedy-accepted reactions register against the actual IRC-determined reactants, which may differ from the seeding context's compounds.

\noindent\begin{tabular}{lr}
\toprule
Parameter & Value \\
\midrule
RMSD match threshold (strict) & \SI{1.0}{\angstrom} \\
Greedy fallback (single + merge) & enabled \\
Max fragments / side (greedy) & 2 \\
\bottomrule
\end{tabular}

The \SI{1.0}{\angstrom} threshold is much looser than the \SI{0.20}{\angstrom} minimum-deduplication cut used at the graph-commit stage because the test here is ``same chemical species, possibly different conformer'' --- conformer-collapse in the per-compound PES graph handles the disambiguation downstream. Greedy mode is the single most consequential pipeline choice: it allows the diffusion model to discover reactions the sampler did not anticipate, but it imports any biases of the model.


%-----------------------------------------------------------------------------
\subsubsection{Stage 7 --- Fragmentation detection}
\label{sec:fragmentation}
%-----------------------------------------------------------------------------

The product side is decomposed into connected components by the covalent-radius-scaled distance criterion (\S\ref{sec:shared-fragments}), with the standard validity rules (non-empty; satisfies the tetrose composition cap; single-atom only if hydrogen). The fragmentation is all-or-nothing: if any component fails validity, the context is dropped --- partial validity would yield inconsistent reaction stoichiometry.

For each valid fragment, the canonical SMILES is computed and an integer charge is inferred from the perceived bond pattern (\S\ref{sec:shared-smiles}); the DFT refinement (\S\ref{sec:kinetics-dft}) re-derives charges from a wavefunction calculation downstream. When a product fragments into multiple components, each fragment is independently re-relaxed on its own PES (FIRE, \S\ref{sec:shared-fire}; $f_\text{max} = \SI{0.005}{eV/\angstrom}$, max $200$ steps) to remove the residual inter-fragment distortion left by the combined-system IRC.

\noindent\begin{tabular}{lr}
\toprule
Parameter & Value \\
\midrule
Covalent-radius scale & 1.3 \\
Per-fragment composition cap & $\{C{:}5,H{:}10,O{:}5\}$ \\
Single-atom rule & H only \\
Per-fragment $f_\text{max}$ & \SI{0.005}{eV/\angstrom} \\
Per-fragment max steps & 200 \\
\bottomrule
\end{tabular}

Per-fragment relaxation matters because the IRC product is a minimum of the combined (interacting) system, not of each fragment in isolation; un-corrected, the per-fragment energies would not be comparable to those of the same fragments encountered in other reactions.


%-----------------------------------------------------------------------------
\subsubsection{Stage 8 --- Backward barrier gate}
\label{sec:bwd-barrier}
%-----------------------------------------------------------------------------

Symmetric to Stage 4. The backward barrier is $\Delta E_\text{bwd} = E_\text{TS} - E_\text{product}$ (summed over fragments if applicable), accepted iff $\Delta E_\text{bwd} \in (\SI{0}{eV}, \SI{2.7211}{eV}]$, with \emph{strict} positivity (a negative backward barrier means a non-stationary or non-converged IRC). The barrier is computed by trajectory force integration (\S\ref{sec:shared-traj}). Surviving contexts pass to the graph-addition step, which canonicalises and registers reactant/product compounds, deduplicates minima against the per-compound PES graphs (three-tier RMSD + energy fast-path, then NEB at \SI{0.05}{eV} barrier threshold; \S\ref{sec:shared-rmsd}, \S\ref{sec:shared-neb}), and either creates a new reaction row or updates an existing lower-barrier variant. Distinct sampling contexts that converged to the same reaction collapse at this point.


%-----------------------------------------------------------------------------
\subsubsection{Pipeline funnel statistics}
\label{sec:funnel}
%-----------------------------------------------------------------------------

\begin{table}
\centering
\begin{tabular}{lrr}
\toprule
Stage & Survivors & Cumulative survival \\
\midrule
Submitted (Stage 1)                                  & $3\,814\,199$         & $100.00$\% \\
After Stages 2--3 (merge prep + diffusion + dedup)   & $3\,814\,199$         & $100.00$\% \\
After Stage 4 (fwd barrier)                          & $1\,798\,768$         & $47.16$\%  \\
After Stages 5--6 (Hessian + IRC + endpoint id)      & $172\,267$            & $4.52$\%   \\
After Stage 7 (fragmentation)                        & $172\,267$            & $4.52$\%   \\
After Stage 8 (bwd barrier)                          & $127\,670$            & $3.35$\%   \\
\midrule
Added to graph (post-dedup)                          & $30\,493$             & $0.80$\%   \\
\bottomrule
\end{tabular}
\caption{Generative-loop funnel for the main run, aggregated over all batches. For comparison, the PES loop (SI-\ref{si:pes}) added $16\,879$ of $69\,450$ escaped reactions, so the generative loop contributed $\sim 64\%$ of the reactions in the graph.}
\label{tab:funnel}
\end{table}

\paragraph{Attrition.}
\textbf{Stage 4} kills $\sim 53\%$ of denoised candidates; almost all rejections are above the upper bound ($> \SI{2}{eV}$ overflow: $1{,}093{,}685$ merge + $1{,}050{,}310$ single), corresponding to strained or partially dissociated denoiser outputs. \textbf{Stages 5--6} kill $> 90\%$ of the remainder: of $1{,}798{,}768$ entering IRC, $418{,}475$ fail the Hessian check; $125{,}090$ and $89{,}898$ fail forward/backward relaxation; $706{,}221$ have both endpoints relax back to the source (\textsc{both\_match}); $412{,}454$ have neither match; only $46{,}630$ pass strict matching. Greedy fallback engages on the $1{,}118{,}675$ strict-mismatched cases and accepts $125{,}637$ (the dominant greedy rejection is identical SMILES multisets on both sides: $961{,}323$). \textbf{Stage 7} rejected zero contexts in the main run --- every IRC-passing geometry fragmented cleanly. \textbf{Stage 8} kills a further $44{,}597$ ($\sim 26\%$) on near-zero backward barriers. \textbf{Graph dedup} removes $97{,}177$ ($\sim 76\%$ of post-Stage-8 survivors), reflecting the fact that many distinct sampling contexts converge to the same reaction.

\paragraph{Merge vs.\ single.}
Submitted: $2{,}613{,}816$ single, $1{,}200{,}383$ merge. After Stage 4: $1{,}706{,}049$ single ($65.3\%$) vs.\ $92{,}719$ merge ($7.7\%$) --- merged complexes are far more often rejected at the cheap energy gate because the fragment-sum baseline is much further from the TS. After Stage 5: $164{,}070$ vs.\ $8{,}197$ ($\sim 9\times$ asymmetric). Greedy-accepted: $120{,}808$ single vs.\ $4{,}829$ merge --- single contexts depend on greedy almost universally, merges mostly succeed under strict matching.
